\documentclass[]{article}

\usepackage{amsmath}
\usepackage{multirow}
\usepackage{natbib}
\usepackage{graphicx} 


\begin{document}

In this study, we investigated the transport of passive Lagrangian tracers in a direct numerical simulation of three-dimensional, subsonic turbulence driven by large-scale solenoidal forcing. The primary goal was to understand how the rotational nature of the energy injection and the resulting coherent vortex structures influence the statistical properties of particle dispersion, specifically its anisotropy and potential for anomalous transport. To this end, we analyzed the trajectories of thousands of tracers, focusing on the Mean-Square Displacement, its decomposition relative to the large-scale flow, the duration of vortex trapping events, and the evolution of displacement probability distributions.

Our analysis revealed two principal findings. First, we identified a persistent and strong anisotropy in the dispersion process. After an initial ballistic phase where particle motion is aligned with the large-scale velocity, the transport becomes dominated by dispersion perpendicular to the local flow direction. The anisotropy ratio stabilizes at a value of approximately 0.52, indicating that transverse transport is nearly twice as effective as parallel transport. This transverse-dominant dispersion is a direct kinematic signature of the solenoidal forcing, where large-scale rotational motions sweep particles more efficiently across the flow than advection carries them along it.

Second, we examined the hypothesis that vortex trapping is a primary driver of anomalous transport. While we confirmed that tracers are transiently captured by coherent vortex structures, we found that these trapping events are remarkably brief. The characteristic residence time within a vortex is only about 7\% of a large-eddy turnover time. These short trapping durations, likely limited by three-dimensional vortex instabilities, are insufficient to generate the long-term memory required for sustained superdiffusion or Lévy-flight statistics. Consequently, the probability distributions of tracer displacements do not develop heavy tails; they are nearly Gaussian at intermediate times and become platykurtic (light-tailed) at late times due to the geometric constraints of the finite periodic domain.

From these results, we conclude that in three-dimensional solenoidal turbulence, the forward energy cascade and inherent vortex dynamics effectively suppress long-range temporal correlations in tracer trajectories. Although the solenoidal forcing imposes a strong and persistent anisotropy on the transport, it does not lead to a breakdown of classical diffusion in the long-time limit. The system robustly returns to a diffusive regime, highlighting that the presence of coherent structures alone is not a sufficient condition for generating anomalous transport.

\end{document}
                